% Experiment 2 (simulation budget), one-paragraph summary. Figure: exp2_mean_reward.pdf (plot_exp2.py).
\paragraph{Experiment 2: simulation budget.}
To measure sample efficiency, we run the same five planners with simulation budgets $N$ from 100 to 2000 and report
the mean reward averaged over $\theta_{\text{exec}}$ on FrozenLake, CliffWalking and Gambler
(Figure~\ref{fig:exp2}). On the two grid worlds, Robust-UCT stays far below all graph planners at every budget and
improves only slowly with $N$: on a tree, each node sees only the simulations that follow its own history, so even a
large budget leaves the worst-case estimate of each node with few samples. The graph planners improve much faster
because every simulation that reaches a state updates the same node, so the samples of different trajectories are
reused. The benefit of this reuse is largest at small budgets and for the most aggressive merging: GSF-Robust-UCT,
which merges states across depths, is already the best planner with the smallest budgets, while GS-Robust-UCT, which
merges only within a depth, needs a larger budget to catch up with it on FrozenLake and stays below it on CliffWalking.
On CliffWalking, the robust graph planners also keep improving with $N$ after the nominal graph planners level off,
since the robust backup needs accurate next-state distributions, which is exactly what pooled samples provide. On Gambler, the average
over $\theta_{\text{exec}}$ is nearly identical for all planners at every budget, because their curves are close to
symmetric around $p_h = 0.5$, so this metric does not separate them there. Overall, sample reuse through the graph
makes the robust planners markedly more sample-efficient than the tree-based Robust-UCT, and full merging gives the
largest gain when the simulation budget is small.
